\originalchapter{相互作用、凝聚与临界点}\label{ch:interactions}
\originalsection{累积量展开与Mayer函数}
理想气体无法解释液化. 设$H=\sum p_i^2/(2m)+U$, $U=\sum_{i<j}v(r_{ij})$, 把动量积分做完得$Z=Z_0\avg{\e^{-\beta U}}_0$, 下标0表示位置均匀独立的理想气体平均. 对数矩母函数的Taylor系数称累积量, 所以
\[
\ln Z=\ln Z_0-\beta\avg U_0+\frac{\beta^2}2\Var_0U-\frac{\beta^3}{6}\avg{U^3}_{0,c}+\cdots.
\]
前两个位置积分在忽略边界、势短程且可积时为$\avg U_0=N(N-1)\int v(r)\dd^3r/(2V)$以及$\Var_0U\simeq N(N-1)\int v(r)^2\dd^3r/(2V)$. 无共同粒子的两键由独立性消去; 只共一个粒子的两键因平移不变平均因子化, 也消去, 剩相同粒子对. 有边界或外场时这些因子化需重新检查.

累积量图中每个粒子坐标为顶点, 每个势因子为键. 不连通片段的乘积在对数中消去, 可由单个顶点拆开并因子化的部分也不贡献固定密度累积量. 联通图只留下一个整体平移积分$V$, 保证$\ln Z$广延. 若$v$在硬核处无穷, $\beta v$展开各阶发散. 将同一对顶点的所有键先求和, 得Mayer函数$f(r)=\e^{-\beta v(r)}-1$, 硬核内为$-1$, 远处趋0, 改善展开参数.

\originalsection{集团展开与维里系数}
定义活度$x=\e^{\beta\mu}/\lambda_T^3$, 巨配分函数$\Xi=\sum_Nx^N\int\prod\dd^3q_i\prod_{i<j}(1+f_{ij})/N!$. 展开乘积就是对全部图求和. $l$个标记粒子的联通图积分总和为$b_l$, $b_1=V$, $b_2=V\int f\dd^3r$. 如果图由$n_l$个$l$粒子联通片段组成, 标签分配数为$N!/\prod_l[n_l!(l!)^{n_l}]$. 对所有$N$求和移除$\sum ln_l=N$约束, 各$n_l$独立指数求和, 因而
\[
\ln\Xi=\sum_{l\ge1}\frac{x^lb_l}{l!},\quad
\beta P=\sum_{l\ge1}\frac{x^l\bar b_l}{l!},\quad
n=\sum_{l\ge1}\frac{x^l\bar b_l}{(l-1)!},\qquad\bar b_l=\lim b_l/V.
\]
这段组合推导解释“对数只含联通图”; 不能把所有联通图都称为不可约图.

逆解$x=n-\bar b_2n^2+(2\bar b_2^2-\bar b_3/2)n^3+\cdots$, 代回压强得
\[
\beta P=n+B_2n^2+B_3n^3+\cdots,\quad
B_2=-\frac12\int f(r)\dd^3r,\quad
B_3=-\frac13\int f_{12}f_{13}f_{23}\dd^3r_{12}\dd^3r_{13}.
\]
$\bar b_3$有三个链图和一个三角图, 消去活度后链图相消, 所以$B_3$只剩三角. 更高阶$B_l$由无割点不可约图组成. 展开适合低密度、稳定短程势; Coulomb等未屏蔽长程势不能假设这些积分有限.

硬球直径$r_0$外有$v(r)=-u_0(r_0/r)^6$, 高温\textbf{$\beta u_0\ll1$}时外部$f\simeq\beta u_0(r_0/r)^6$. 核内积分为$-4\pi r_0^3/3$, 外部为$4\pi\beta u_0r_0^3/3$, 故$B_2=b(1-\beta u_0)$, $b=2\pi r_0^3/3$. 原课PDF在该近似处的“$\beta u_0\gg1$”与其Taylor展开不相容, 本稿依据实际展开条件改为$\ll1$, 这是明示校正.

\originalsection{van der Waals平均场与稳定性}
低密度排除体积的逐粒子积分近似为$\prod_{j=0}^{N-1}(V-j\Omega)\simeq(V-N\Omega/2)^N$, $\Omega=4\pi r_0^3/3$, 只在相应低密度精度上使用. 均匀吸引能$-uN^2/(2V)$使
\[
Z\simeq\frac{(V-Nb)^N}{N!\lambda_T^{3N}}\e^{\beta aN^2/V},
\quad P=\frac{n\kb T}{1-bn}-an^2,\quad a=u/2.
\]
吸引降低压强, 排斥提高压强. 正则系综中动量分布与位置分布分离, 吸引使壁边密度相对体内降低; 沿某条确定轨道做动力学解释时, 粒子接近壁面也会改变速度和停留时间. 两种记账必须一致, 不重复把两个解释的压强修正相加.

稳定均匀相要求$(\partial P/\partial n)_T\ge0$, 等价$\kappa_T\ge0$. van der Waals低温等温线的负斜率密度区对应负压缩率, 表明均匀近似不稳定. 两相共存须同时满足相同$T,P,\mu$. 用$d\mu=v\dd P$沿等温线积分, 或对自由能取公切线, 得Maxwell等面积式$\int_{v_l}^{v_g}[P(v,T)-P_{\rm coex}]\dd v=0$. 真实共存段是水平线, 不是保留不稳定回环作为可实验准静态路径.

\originalsection{变分估计与热力学极限}
可解参考Hamilton量$H_0$的正则密度$p_0$下, 对经典$Z/Z_0=\avg{\e^{-\beta(H-H_0)}}_0$使用Jensen不等式, 得$F\le F_0+\avg{H-H_0}_0$. 用含参数的$H_0$最小化右端给自由能上界. 对量子不对易Hamilton量, 不能直接把两个矩阵指数拆乘; 同一个Gibbs--Bogoliubov界可由量子相对熵非负性得到, 见第\ref{ch:density}章. 硬核参考体系加平均吸引与密度变分, 产生前述平均场近似, 但其参考硬核自由能本身也若被近似, 就不能继续声称对真实体系有严格上界.

巨正则大体积和可写作$\Xi\sim\sum_n\exp\{\beta V[\mu n-f(n,T)]\}$. 鞍点条件$f_n=\mu$, 稳定相使$\mu n-f$最大. 两个局部最大相等时交换占优, 热力学极限$V^{-1}\ln(\e^{VA}+\e^{VB})\to\max(A,B)$出现不可微. 对有限自由度、稳定势和收敛的有限温配分函数, 原积分通常光滑; 严格相变的奇点需取适当无限系统极限. 不能忽略配分函数收敛失败等例外后断言所有有限系统都绝无奇点.

\originalsection{对应态与临界指数}
用$v=1/n$, 同时解$P_v=P_{vv}=0$, 得$v_c=3b$, $\kb T_c=8a/(27b)$, $P_c=a/(27b^2)$. 约化变量$P_r=P/P_c$, $v_r=v/v_c$, $T_r=T/T_c$满足$P_r=8T_r/(3v_r-1)-3/v_r^2$. 本模型预测$P_cv_c/(\kb T_c)=3/8$, 实际简单气体的比值并非统一3/8, 所以对应态近似不能当成所有液气体系的精确定律.

临界附近, 假设等温线解析, 记$t=T-T_c$, $x=v-v_c$, 可写$P-P_c=A t-a_1tx-c_1x^3/6+\cdots$, $a_1,c_1>0$. 沿临界等容线$\kappa_T\sim t^{-1}$, 临界等温线$P-P_c\sim-x^3$, 低于临界的共存体积差由$a_1|t|x\sim c_1x^3/6$给$|x|\sim|t|^{1/2}$. 这些是平均场指数$\gamma=1,\delta=3,\beta_{crit}=1/2$; 实际三维简单流体的临界指数约$1.24,4.8,0.326$, 原课以粗略$1.3,5.0,0.3$说明偏离. 符号$\beta_{crit}$与逆温度$\beta$含义不同. 接近临界时强涨落破坏解析平均场假设, 普适性与重整化属于未选的8.334后续课程, 本稿不冒称在此已证明这些真实指数.
